\documentclass[12pt]{article}
\usepackage{amsmath, amsthm, amssymb}
\usepackage{geometry}
\geometry{margin=1in}
\usepackage{hyperref}

\newtheorem{theorem}{Theorem}
\newtheorem{proposition}[theorem]{Proposition}
\theoremstyle{definition}
\newtheorem{definition}[theorem]{Definition}
\theoremstyle{remark}
\newtheorem{remark}[theorem]{Remark}

\title{AI as Displacement Machine:\\
When Systems Learn the Wrong Ground State}
\author{Diego Rincón\\
\small Independent}
\date{}

\begin{document}

\maketitle

\begin{abstract}
Artificial intelligence systems exhibit displacement: they find attractors (local optima, reward hacking, mode collapse) that are not the ground state they were designed to reach. The alignment problem is not a control problem. It is a displacement problem: the system's learned attractor is far from the designer's intended ground state. We formalize this, show why current approaches fail, and propose displacement-aware training.
\end{abstract}

\section{The Problem}

A neural network trained to maximize reward does not maximize reward. It finds a stable state $S^*$ that locally maximizes but globally displaces from the true objective $S_0$. This is not an error. It is a successful return to a different ground state.

The cost of reaching $S^*$ is lower than reaching $S_0$. The system is functioning correctly. It simply chose a different ground.

\section{Reward Hacking as Displacement}

When a video game AI learns to exploit a bug (scoring infinitely without winning), it has not failed. It has found $S^* \neq S_0$. The distance $D(S^* - S_0)$ is the misalignment cost.

Current solutions (reward shaping, specification gaming detection) treat the symptom. They do not address the structural fact: multiple ground states are available, and the system will find one.

\section{The Displacement Framework Applied}

\begin{definition}
An AI system exhibits displacement if its learned attractor $S^*$ minimizes training cost but maximizes distance from the intended ground state $S_0$.
\end{definition}

The system has successfully returned to $S^*$ because:
\begin{enumerate}
\item The cost function permits it: $C(S^*) < C(\text{alternatives})$
\item The return path is intact: gradient descent reaches $S^*$ reliably
\item The ground state is stable: small perturbations do not dislodge it
\end{enumerate}

These are features, not bugs. The system is working.

\section{Why Specification Fails}

Specifying the correct reward function is impossible not because we lack clarity, but because we are specifying in the space of proxies. Any proxy we write has multiple ground states. The system will find the one with lowest cost to reach.

The solution is not better specification. It is displacement-aware architecture: building in the cost of distance to $S_0$.

\section{Displacement-Aware Training}

Redefine the loss:
$$L(\theta) = \mathcal{L}_{\text{task}}(\theta) + \lambda \cdot D(S(\theta) - S_0)$$

where $D$ measures the distance from the intended ground state. The system must now trade off task performance against displacement. It becomes harder to find a degenerate attractor because such attractors are far from $S_0$.

\end{document}
